% Generated draft: requires mathematical and source review.
\documentclass[11pt]{article}
\usepackage[margin=25mm]{geometry}
\usepackage{fontspec}
\setmainfont{Latin Modern Roman}
\usepackage{amsmath,amssymb,mathtools,unicode-math}
\setmathfont{Latin Modern Math}
\usepackage{enumitem,needspace,xurl,hyperref}
\setlength{\parindent}{0pt}
\setlength{\parskip}{5pt}
\setlength{\emergencystretch}{3em}
\newcommand{\N}{\mathbb{N}}
\newcommand{\Z}{\mathbb{Z}}
\newcommand{\Q}{\mathbb{Q}}
\newcommand{\R}{\mathbb{R}}
\newcommand{\C}{\mathbb{C}}
\newcommand{\F}{\mathbb{F}}
\newcommand{\A}{\mathbb{A}}
\newcommand{\PP}{\mathbb{P}}
\newcommand{\E}{\mathbb{E}}
\newcommand{\eps}{\varepsilon}
\newcommand{\dd}{\,\mathrm{d}}
\begin{document}
\section*{UnsolvedMath ID: 30004534}
\textbf{Source problem code:} catalogue entry
\subsection{Definitions and mathematical statement}
The supplied record does not specify the coefficient field, ambient rings, grading, transition maps, or precise meaning of a \(\operatorname{Sym}\)-invariant chain. Consequently, it does not determine a complete class of chains over which the question is quantified. The following gives a conditional formulation of its Lefschetz assertion.

Suppose that a fixed field \(K\) and homogeneous ideals \(I_n\) in graded commutative \(K\)-algebras \(R_n\) are given, and that each quotient
\[
A_n=R_n/I_n=\bigoplus_{i\geq 0}(A_n)_i
\]
is a finite-dimensional standard graded \(K\)-algebra: \((A_n)_0=K\), and \(A_n\) is generated by \((A_n)_1\). These are additional hypotheses for defining the usual graded Lefschetz properties; they are not fully specified in the source record. Saying that \(I_n\) has a Lefschetz property is interpreted here as saying that \(A_n\) has that property.

A finite-dimensional standard graded \(K\)-algebra \(A\) has the \emph{weak Lefschetz property}, denoted \(\operatorname{WLP}(A)\), if there exists \(\ell\in A_1\) such that, for every integer \(i\geq0\), multiplication induces a map
\[
\times\ell:A_i\longrightarrow A_{i+1}
\]
of rank \(\min\{\dim_K A_i,\dim_K A_{i+1}\}\). It has the \emph{strong Lefschetz property}, denoted \(\operatorname{SLP}(A)\), if there exists \(\ell\in A_1\) such that, for all integers \(i,d\geq0\), the map
\[
\times\ell^d:A_i\longrightarrow A_{i+d}
\]
has rank \(\min\{\dim_K A_i,\dim_K A_{i+d}\}\). These definitions require the Lefschetz element to be defined over \(K\); whether the intended question instead permits extending the field is unspecified.

\emph{Proposed eventual-constancy question.} For every chain satisfying the intended, but incompletely specified, \(\operatorname{Sym}\)-invariance assumptions, is it true that, separately for each \(P\in\{\operatorname{WLP},\operatorname{SLP}\}\),
\[
\bigl(\exists N\geq1\;\forall n\geq N:\ P(A_n)\bigr)
\quad\lor\quad
\bigl(\exists N\geq1\;\forall n\geq N:\ \neg P(A_n)\bigr)?
\]
The threshold and the eventual truth value may depend on the chain and on the chosen property. This separate interpretation follows the supplied cleaned statement. A precise unconditional formulation requires the missing algebraic and symmetry hypotheses. Moreover, the supplied literature assessment says that the extracted report passage is only a general remark about possible limit behavior and identifies no definite theorem or conjecture; attribution of this precise question to that passage therefore remains unresolved.
\subsection{Short English statement}
The proposed question asks whether the weak and strong Lefschetz properties each eventually stabilize along every intended symmetric chain of Artinian quotients. Each property may eventually hold everywhere or eventually fail everywhere, with its own threshold. The source leaves the class of chains and essential algebraic hypotheses unspecified, and its dated assessment questions whether the report contains this precise conjecture.
\subsection{Sources}
\begin{itemize}
\item Local UnsolvedMath registry, \texttt{data/unsolvedmath/problems.json}: permanent ID \(30004534\), source problem code \texttt{catalogue entry}, title ``Eventual Lefschetz Dichotomies for Symmetric Ideal Chains.'' This draft uses the supplied complete registry record; its registry location has not been independently inspected.
\item Supplied report citation: \emph{Lefschetz Properties in Algebra, Geometry and Combinatorics} (2021), Oberwolfach Reports, pp.\ 1539--1600, DOI \texttt{10.4171/owr/2020/31}, \url{https://doi.org/10.4171/owr/2020/31}. This same URL is supplied under literature sources with the label ``OWR report.'' No individual proposer or pinpoint location within the report is supplied. The record gives proposed year \(2020\).
\item Supplied literature assessment, checked on \texttt{2026-08-22}: the record labels the problem open, while stating that the extracted passage contains only a general remark about possible limit behavior and no identifiable theorem or conjecture. This is supplied dated curation, not an independently verified literature assessment. The supplied background explicitly calls for independent verification of time-sensitive status and cited claims.
\item The registry labels the statement \texttt{corrected\_verified} and says its mathematical content was preserved while presentation and omitted definitions were repaired. This is supplied metadata, not independent mathematical or source verification; it does not resolve the qualifications above.
\end{itemize}
\end{document}
